\documentclass[12pt]{article}
\usepackage[margin=1in]{geometry}
\usepackage{setspace}
\usepackage{lmodern}
\renewcommand{\labelitemi}{\ensuremath{\bullet}}

\begin{document}
\doublespacing

\begin{center}
\textbf{Puntos Destacados de la Investigación} \\
\vspace{1em}
\textbf{Identifying Rational Types in Unknown Environments under an Indirect Dynamic Mechanism} \\
Humberto Bernal
\end{center}

\vspace{2em}

\begin{itemize}
    \item Modela interacciones estratégicas entre el Estado e insurgencias (FARC, ELN, AUC) mediante cribado dinámico indirecto.
    \item El proceso Mano de Dios--Guadalupe introduce ruido de máxima entropía para sustentar la inferencia bayesiana bajo fricción de conflicto.
    \item El Teorema de Identificación de Tipos Racionales demuestra convergencia casi segura hacia el tipo insurgente ideológico o criminal.
    \item Riesgos en competencia microeconométricos calibran la dinámica empírica del secuestro en cuatro décadas de conflicto colombiano.
    \item La evaluación de política contraterrorista cuantifica los trade-offs estructurales entre interdicción financiera ($\alpha$) y presión táctica ($\gamma$).
\end{itemize}

\end{document}
